\chapter{Desenvolvimento}
\label{cap:desenvolvimento}

Este capítulo descreve como o Laboratório 404 foi construído: a arquitetura
geral, a especificação que organizou o trabalho e o desenvolvimento de cada
prova de conceito, com as decisões de engenharia relevantes e os defeitos
encontrados e corrigidos. O capítulo segue a ordem cronológica do projeto,
que é também a ordem de dependência entre os subsistemas.

\section{Visão geral da arquitetura}
\label{sec:arquitetura}

O sistema é composto por quatro blocos: o \textbf{jogo} (motor Godot, com
\emph{gameplay}, física, interface, áudio e persistência); os \textbf{assets
3D} (modelados no Blender e exportados em GLB); o \textbf{adaptador de IA}
(serviço local em Python que valida e media a conversa); e a \textbf{IA}
propriamente dita (modelo de linguagem executado localmente pelo Ollama). A
Figura~\ref{fig:arquitetura} apresenta o conjunto.

\begin{figure}[htbp]
\centering
\includegraphics[width=0.95\textwidth]{01_arquitetura.png}
\caption{Arquitetura do Laboratório 404: o jogo consome assets do Blender e ações de entrada; a IA vive em processos separados, atrás de um adaptador que valida entrada e saída.}
\label{fig:arquitetura}
\end{figure}

Duas decisões estruturais merecem destaque. A primeira é manter a IA
\emph{fora} do processo do jogo: se o adaptador ou o modelo falham, o jogo
continua jogável e responde com textos de \emph{fallback}. A segunda é tratar
o adaptador como fronteira de segurança: é nele que a entrada é validada, que
o contexto do jogo é injetado no \emph{prompt} e que a saída do modelo é
reduzida a texto e a uma intenção pertencente a um conjunto finito.

\section{Especificação e organização incremental}
\label{sec:especificacao}

Antes da implementação, o projeto foi especificado com requisitos funcionais
(FR) e não funcionais (NFR), critérios de aceitação no formato
DADO/QUANDO/ENTÃO e uma matriz de rastreabilidade que associa cada requisito
ao teste correspondente. Os requisitos foram agrupados por prova de conceito,
e cada grupo tinha um critério de avanço explícito: a etapa só era considerada
concluída com a suíte verde e as evidências registradas. O
\autoref{ap:requisitos} reproduz a especificação consolidada; o
\autoref{ap:testes} detalha a suíte.

\section{\poc{1} --- Fundação jogável}
\label{sec:poc1}

O primeiro incremento estabeleceu o esqueleto do jogo: um jogador que se
movimenta, uma câmera, um HUD mínimo, menu de pausa e --- decisão que se
provou estruturante --- um mapa de entrada inteiramente baseado em ações
semânticas.

\subsection{Entrada semântica e diagnóstico de controles}

O jogo consome ações nomeadas (por exemplo, \texttt{move\_forward},
\texttt{look\_right}, \texttt{interact}, \texttt{pause}) e nunca índices
numéricos de botão ou eixo. A razão é prática: adaptadores USB convertem
controles de maneiras distintas, e um índice fixo deixa de funcionar quando o
adaptador muda. Com ações semânticas, o mapeamento físico fica isolado na
configuração do motor. O Quadro~\ref{quad:input} ilustra o padrão usado em
todo o \emph{gameplay}.

\begin{quadro}[htbp]
\centering
\caption{Padrão de leitura de entrada usado no gameplay (GDScript).}
\label{quad:input}
\begin{verbatim}
var direcao := Input.get_vector(
    "move_left", "move_right",
    "move_forward", "move_backward")
# NUNCA usar indices crus, como Input.is_joy_button_pressed(0, 1)
\end{verbatim}
\par\vspace{2pt}{\footnotesize Fonte: elaborado pelo autor (2026), a partir do código do projeto.}
\end{quadro}

Para calibrar adaptadores desconhecidos, o projeto inclui uma tela de
diagnóstico de entrada que lista os dispositivos conectados, os eixos e os
botões em tempo real, com zona morta configurável (padrão de 0,15) e
tratamento de conexão e desconexão a quente (\emph{hotplug}) sem interromper a
execução. Essa tela é a única parte do projeto autorizada a ler índices
numéricos; o \emph{gameplay}, jamais.

\subsection{Jogador, câmera, HUD e pausa}

O jogador é um corpo com física que respeita gravidade e colisão. A câmera é
controlada por mouse e, quando o adaptador expõe os eixos do analógico
direito, por este. O HUD exibe o objetivo corrente e o convite de interação
contextual; o menu de pausa (\texttt{Esc}/\emph{Start}) dá acesso ao
diagnóstico de entrada, ao salvamento e à saída. Essas funcionalidades
correspondem aos requisitos FR-002 a FR-006 e foram verificadas
automaticamente (critérios CA-001 a CA-005), com duas ressalvas manuais: o
uso do mouse com captura de cursor e a operação física do joystick.

\section{\poc{2} --- Laboratório funcional}
\label{sec:poc2}

O segundo incremento transformou o esqueleto em um laboratório: uma sala
modelada no Blender, interação por proximidade, inventário, missão e estados
de equipamento.

\subsection{Pipeline de assets: do Blender ao jogo}
\label{sec:pipeline_assets}

Os assets 3D seguem convenções fixas que tornaram a integração previsível:
1 unidade do Blender equivale a 1 metro; as transformações são aplicadas antes
da exportação; os objetos são nomeados por função (Quadro~\ref{quad:prefixos});
e a colisão é embutida no próprio asset por meio do sufixo \texttt{-col} ---
o importador do motor converte essas malhas em corpos estáticos, mantendo a
malha visível original apenas para renderização.

\begin{quadro}[htbp]
\centering
\caption{Prefixos de nomenclatura adotados na modelagem.}
\label{quad:prefixos}
\footnotesize
\begin{tabular}{@{}>{\raggedright\arraybackslash}p{2.0cm} >{\raggedright\arraybackslash}p{11.0cm}@{}}
\toprule
Prefixo & Uso \\
\midrule
\texttt{ENV\_} & arquitetura (paredes, piso, teto, estruturas) \\
\texttt{PROP\_} & equipamento ou objeto cenográfico \\
\texttt{INT\_} & objeto interativo \\
\texttt{COL\_} & colisor auxiliar \\
\texttt{FX\_} & efeito visual \\
\texttt{NPC\_} & personagem \\
\texttt{DEC\_} & decoração \\
\bottomrule
\end{tabular}
\par\vspace{2pt}{\footnotesize Fonte: elaborado pelo autor (2026), a partir das convenções do projeto.}
\end{quadro}

A exportação é feita por script, com parâmetros fixos: formato binário GLB,
transformações aplicadas, eixo vertical convertido, e \emph{exclusão} de
câmeras e luzes do arquivo (a iluminação é reconstruída no motor, o que
permite ajustá-la sem reexportar). O Quadro~\ref{quad:export} reproduz a
chamada de exportação de referência. Depois de cada exportação, o projeto
executa a importação de recursos por linha de comando e a suíte de testes,
garantindo que o asset novo não quebrou nenhuma verificação.

\begin{quadro}[htbp]
\centering
\caption{Exportação de referência do asset 3D (script do Blender).}
\label{quad:export}
\begin{verbatim}
bpy.ops.export_scene.gltf(
    filepath="assets/lab404/lab_sala_poc2.glb",
    export_format='GLB', export_apply=True, export_yup=True,
    export_cameras=False, export_lights=False)
\end{verbatim}
\par\vspace{2pt}{\footnotesize Fonte: elaborado pelo autor (2026), a partir do pipeline do projeto.}
\end{quadro}

O asset final da sala contém 342 nós, 37 materiais e 35 nós de colisão,
totalizando 1.018,7~KB --- números obtidos por inspeção programática do
arquivo, não por estimativa. O emissivo dos materiais usa a extensão
\texttt{KHR\_materials\_emissive\_strength} do glTF, o que permitiu acender
telas e sinalizações com intensidade controlada, sem estourar para branco no
renderizador compatível.

\subsection{A sala e o detalhamento visual}
\label{sec:sala}

A sala do laboratório foi modelada a partir de uma referência visual de
ambiente industrial e refinada iterativamente. A
Figura~\ref{fig:referencia} mostra a referência de arte usada e a
Figura~\ref{fig:sala_geral} mostra o resultado em jogo, no estado inicial ---
energia isolada, iluminação de emergência e a missão pendente no HUD.

\begin{figure}[htbp]
\centering
\includegraphics[width=0.8\textwidth]{referencia_arte.png}
\caption{Referência visual utilizada para a reconstrução da sala (imagem de referência adotada no projeto).}
\label{fig:referencia}
\end{figure}

\begin{figure}[htbp]
\centering
\includegraphics[width=0.8\textwidth]{sala_geral.png}
\caption{Sala do Laboratório 404 no estado inicial: energia isolada, iluminação de emergência, técnico (NPC) à esquerda e bomba em falha. Captura do jogo em 1280$\times$720.}
\label{fig:sala_geral}
\end{figure}

O detalhamento seguiu um princípio econômico: cada elemento adicional precisava
ter função visual ou narrativa. O CLP ganhou ranhuras de ventilação e quatro
LEDs de status que piscam em períodos distintos; o painel de encaixe recebeu
faixas zebradas de perigo, triângulo de alerta elétrico, esquema de fiação e
molduras dos slots; a bomba recebeu braçadeiras, apoios ancorados ao piso,
longarinas com parafusos sextavados, manômetro analógico com ponteiro e
válvulas de alívio; e a bancada de trabalho recebeu luminária própria. As
Figuras~\ref{fig:clp} a \ref{fig:bancada} ilustram esse conjunto.

\begin{figure}[htbp]
\centering
\includegraphics[width=0.72\textwidth]{clp.png}
\caption{Estação do CLP, com ranhuras de ventilação, LEDs de status e diagrama \emph{ladder} na tela.}
\label{fig:clp}
\end{figure}

\begin{figure}[htbp]
\centering
\includegraphics[width=0.72\textwidth]{painel.png}
\caption{Painel de encaixe, com sinalização de perigo, esquema de fiação e slots para o fusível e o módulo.}
\label{fig:painel}
\end{figure}

\begin{figure}[htbp]
\centering
\includegraphics[width=0.72\textwidth]{bomba.png}
\caption{Conjunto da bomba hidráulica, com manômetro, tubulações e ancoragens mecânicas, em estado de falha.}
\label{fig:bomba}
\end{figure}

\begin{figure}[htbp]
\centering
\includegraphics[width=0.72\textwidth]{bancada.png}
\caption{Bancada de trabalho com luminária dedicada, onde está o item coletável inicial.}
\label{fig:bancada}
\end{figure}

Duas decisões técnicas do detalhamento merecem registro porque contrariam
expectativas comuns de quem vem da modelagem tradicional: LEDs que piscam
\emph{não} podem ser animados no glTF (o formato anima apenas translação,
rotação e escala, não emissão de material --- de modo que o piscar é feito por
script no motor, alternando a visibilidade dos nós dos LEDs); e textos de
tela são implementados como nós de texto do próprio motor, sobre a tela
emissiva do modelo, evitando dependência de fontes dentro do GLB e
permitindo ajustes de idioma sem reexportar o asset.

\subsection{Interação, inventário e missão}
\label{sec:interacao}

A interação é um sistema desacoplado: um raio de proximidade detecta objetos
interativos, exibe o convite contextual no HUD e, quando a ação de interagir é
acionada, emite um sinal que o objeto trata. O inventário é diegético (um kit
técnico), em vez de uma grade abstrata: o jogador coleta o fusível F-17, o
módulo RS-404 e uma chave de manutenção, e os instala nos pontos
correspondentes da sala.

A missão principal, \texttt{RESTAURAR\_COMUNICACAO}, encadeia seis passos
dependentes (Figura~\ref{fig:missao}): localizar o fusível, localizar o
módulo, instalar ambos no painel, reiniciar o CLP, diagnosticar a bomba e
abrir a porta do corredor. Cada passo só é aceito quando o anterior está
concluído, e o HUD reflete o avanço. O encadeamento é o que dá ao jogo sua
estrutura didática: o jogador percorre uma sequência de diagnóstico
semelhante à de um atendimento real.

\begin{figure}[htbp]
\centering
\includegraphics[width=0.95\textwidth]{04_missao.png}
\caption{Missão principal em seis passos dependentes e ordenados.}
\label{fig:missao}
\end{figure}

\subsection{Estados de equipamento}
\label{sec:estados}

Os equipamentos possuem uma máquina de estados explícita --- desligado,
falha, pronto e em operação --- serializável e independente da interface
(Figura~\ref{fig:estados}). A escolha por um estado explícito, em vez de
sinalizações dispersas, permitiu dois efeitos diretos: o mundo (luzes, NPC,
tom da IA) reage ao estado sem conhecer detalhes de implementação, e o
progresso do jogador pode ser salvo e restaurado com fidelidade.

\begin{figure}[htbp]
\centering
\includegraphics[width=0.72\textwidth]{03_estados_equipamento.png}
\caption{Máquina de estados dos equipamentos, serializável e independente da interface.}
\label{fig:estados}
\end{figure}

\subsection{Porta B1: um defeito visual e sua correção}
\label{sec:porta}

A porta do corredor --- obstáculo final da missão --- apresentou um defeito
instrutivo. No asset original, a folha da porta era composta por várias malhas
soltas; apenas a malha de colisão estava vinculada ao mecanismo de abertura.
O resultado visível era um paradoxo: a porta abria logicamente (o jogador
passava), mas o visual da folha permanecia fechado no vão. A correção foi
feita em duas frentes: no Blender, as malhas decorativas (reforço, visor e
placa) foram tornadas filhas do objeto da porta, e no jogo o mecanismo passou
a elevar a folha inteira em 2,15~m e a ocultá-la ao final do curso. As
Figuras~\ref{fig:porta_fechada} e \ref{fig:porta_aberta} mostram o antes e o
depois; os testes ganharam verificações específicas (a folha deve conter os
três detalhes como filhos e subir mais de 1,5~m).

\begin{figure}[htbp]
\centering
\includegraphics[width=0.72\textwidth]{porta_fechada.png}
\caption{Porta B1 fechada, antes de a missão ser concluída.}
\label{fig:porta_fechada}
\end{figure}

\begin{figure}[htbp]
\centering
\includegraphics[width=0.72\textwidth]{porta_aberta.png}
\caption{Porta B1 aberta após a missão: a folha sobe e libera o vão do corredor.}
\label{fig:porta_aberta}
\end{figure}

Dois detalhes adicionais merecem registro no detalhamento dos equipamentos: o
esquema de fiação serigrafado do painel foi reposicionado na margem direita
(na posição original, ficava escondido atrás da tela) e o pato de borracha da
bancada --- um elemento decorativo --- precisou ser elevado acima da
superfície de colisão da bancada, que é mais alta que a superfície visível.
As Figuras~\ref{fig:det_painel} a \ref{fig:det_bomba} ilustram o nível de
detalhe alcançado.

\begin{figure}[htbp]
\centering
\includegraphics[width=0.72\textwidth]{detalhe_painel.png}
\caption{Detalhe do painel: faixas de perigo, triângulo de alerta, esquema de fiação e molduras dos slots.}
\label{fig:det_painel}
\end{figure}

\begin{figure}[htbp]
\centering
\includegraphics[width=0.72\textwidth]{detalhe_telas.png}
\caption{Detalhe das telas: conteúdos em português exibidos sobre as superfícies emissivas.}
\label{fig:det_telas}
\end{figure}

\begin{figure}[htbp]
\centering
\includegraphics[width=0.72\textwidth]{detalhe_bomba.png}
\caption{Detalhe da bomba: braçadeiras, apoios, parafusos sextavados e manômetro.}
\label{fig:det_bomba}
\end{figure}

\section{\poc{3} --- \aria{}: IA local com intenções restritas}
\label{sec:poc3}

O terceiro incremento adicionou a personagem de IA. É o subsistema em que as
decisões de segurança mais importam, e por isso ele é descrito com detalhe.

\subsection{Fluxo de uma mensagem}
\label{sec:fluxo_mensagem}

A Figura~\ref{fig:fluxo} mostra o caminho de uma mensagem. O jogador digita
no terminal do jogo; o cliente (em GDScript) envia um JSON para o adaptador
local; o adaptador valida a entrada, monta o \emph{prompt} com a
personalidade versionada e o estado do jogo, chama o modelo e valida a
resposta antes de devolvê-la. Cada elo da cadeia pode falhar de forma segura:
o pior caso é um texto de \emph{fallback}, nunca um travamento ou uma ação não
prevista.

\begin{figure}[htbp]
\centering
\includegraphics[width=0.95\textwidth]{02_fluxo_conversa.png}
\caption{Cadeia completa de uma mensagem até a resposta: validação em duas camadas e resposta de fallback em caso de falha.}
\label{fig:fluxo}
\end{figure}

\subsection{Contrato de dados}
\label{sec:contrato}

O adaptador expõe um único ponto de conversa e um ponto de verificação de
saúde. A entrada exige mensagem não vazia de até 2000 caracteres e um
dicionário opcional com o estado do jogo; a saída é sempre um objeto com
texto, intenção, alvo opcional, confiança e erro opcional. O
Quadro~\ref{quad:contrato} reproduz o contrato essencial.

\begin{quadro}[htbp]
\centering
\caption{Contrato de dados do adaptador de IA (resumo).}
\label{quad:contrato}
\begin{verbatim}
POST /chat
  request:  {"message": "A bomba nao liga.",
             "game_state": {"bomba": "FAULT", "energia": 62}}
  response: {"text": "...", "intent": "DIAGNOSE",
             "target": "PROP_Bomba", "confidence": 0.94}
validacao: 422 para entrada invalida;
           falha/timeout do modelo => texto de fallback, intent NONE
\end{verbatim}
\par\vspace{2pt}{\footnotesize Fonte: elaborado pelo autor (2026), a partir do adaptador do projeto.}
\end{quadro}

A validação de entrada usa modelos de dados declarativos (Pydantic), o que
garante que nenhuma mensagem fora dos limites chega ao modelo. A validação de
saída é a parte mais importante do desenho: o texto gerado é aceito, mas a
intenção só sobrevive se pertencer ao conjunto permitido --- caso contrário,
é convertida para \texttt{NONE}. A mesma sanitização é repetida no cliente do
jogo, de modo que as duas pontas precisam concordar. O
Quadro~\ref{quad:intents} reproduz o conjunto e a regra.

\begin{quadro}[htbp]
\centering
\caption{Conjunto de intenções permitidas e regra de sanitização (adaptador).}
\label{quad:intents}
\begin{verbatim}
ALLOWED_INTENTS = {"HINT", "DIAGNOSE", "LORE", "STATUS", "NONE"}

def sanitize_intent(valor):
    # remove ruido, normaliza e converte para maiusculas
    candidato = re.sub(r"[^A-Za-z]", "", str(valor or "")).upper()
    return candidato if candidato in ALLOWED_INTENTS else "NONE"
\end{verbatim}
\par\vspace{2pt}{\footnotesize Fonte: elaborado pelo autor (2026), a partir do código do adaptador.}
\end{quadro}

\subsection{Personalidade versionada}
\label{sec:personalidade}

A personalidade da IA não está espalhada pelo código: ela vive em um arquivo
de configuração versionado, com identidade, regras de conduta, tons de voz
(neutra, impaciente, preocupada), saudação e texto de \emph{fallback}. Duas
regras desse arquivo materializam o desenho de segurança adotado: a IA
\emph{não executa comandos, não controla portas, energia nem física}; e nunca
inventa componentes, códigos ou procedimentos inexistentes no laboratório.
O tom pode ser alterado pelo estado do jogo --- por exemplo, quando o jogador
ignora um alerta, a IA responde de forma mais curta e seca ---, o que dá
coerência afetiva à ficção sem ampliar seu poder de ação.

\subsection{Limites de tempo e resposta de segurança}
\label{sec:fallback}

O adaptador aguarda o modelo por até 30 segundos; o cliente no jogo espera um
pouco mais (35 segundos) para acomodar a variabilidade de um modelo local.
Esgotado o prazo --- ou ocorrendo qualquer erro de rede, de validação ou de
modelo ---, a resposta devolvida é o texto de \emph{fallback} com intenção
\texttt{NONE}. Essa é a garantia que sustenta o requisito CA-009: o jogo
nunca trava por causa da IA.

\subsection{No jogo: terminal e estação}
\label{sec:terminal}

No jogo, a conversa acontece em um terminal com histórico de mensagens e
entrada de texto, acessível ao interagir com a estação da IA
(Figura~\ref{fig:aria_estacao}). A estação também exibe uma leitura de status
própria, com identificação da versão do núcleo e do canal de comunicação.

\begin{figure}[htbp]
\centering
\includegraphics[width=0.72\textwidth]{aria_estacao.png}
\caption{Estação física da \aria{} na sala, com tela de status e ponto de interação.}
\label{fig:aria_estacao}
\end{figure}

\section{\poc{4} --- Mundo reativo}
\label{sec:poc4}

O quarto incremento fez o laboratório reagir ao estado do jogo, com quatro
mecanismos.

\subsection{NPC e diálogo contextual}

Um técnico (NPC) comenta o estado do laboratório --- por exemplo, a situação
da bomba --- de acordo com o momento da missão. O personagem foi modelado como
um corpo único de 882 vértices e 10 materiais (macacão, colete refletivo,
luvas, botas, capacete com lanterna e rosto), com a mesma malha servindo de
colisão; a Figura~\ref{fig:tecnico} mostra o resultado em jogo. A fala do NPC
usa o mesmo canal de mensagens do HUD que os avisos do sistema.

\begin{figure}[htbp]
\centering
\includegraphics[width=0.72\textwidth]{tecnico.png}
\caption{Técnico (NPC) na sala, com uniforme, capacete e postura de trabalho.}
\label{fig:tecnico}
\end{figure}

\subsection{Memória em camadas}
\label{sec:memoria}

A memória do jogo foi separada em quatro camadas com ciclos de vida
diferentes (Figura~\ref{fig:memoria}): a sessão atual; o estado persistente
(salvo em arquivo); o conhecimento narrativo acumulado (fatos que o jogador
aprendeu sobre a instalação e a IA); e o histórico conversacional. A separação
evita o problema clássico de misturar \emph{flags} de curto prazo com fatos
que precisam sobreviver ao salvamento --- e permite que a IA receba como
contexto exatamente a camada adequada à pergunta.

\begin{figure}[htbp]
\centering
\includegraphics[width=0.95\textwidth]{07_memoria.png}
\caption{As quatro camadas de memória e estado, com ciclos de vida diferentes.}
\label{fig:memoria}
\end{figure}

\subsection{Eventos e iluminação em camadas}

Os eventos do mundo reagem ao estado: luzes mudam conforme a energia, portas
respeitam permissões e o tom da IA muda quando um alerta é ignorado. A
iluminação foi organizada em duas camadas: a luz geral, com dois estados
(claro e emergência), alinhada às luminárias efetivamente visíveis no modelo
3D; e luzes de acento por zona (emergência, IA, painel, bomba, sala de
controle), que preservam a identidade cromática de cada área mesmo no escuro.
Pouco depois de a energia ser restabelecida, uma câmera de segurança da sala
passa a acompanhar o jogador com um movimento limitado --- uma recompensa
visual pelo avanço na missão.

\section{\poc{5} --- Laboratório procedural}
\label{sec:poc5}

O quinto incremento estendeu o jogo para além da sala modelada à mão. Um
gerador cria uma ala de laboratório a partir de módulos tipados (corredor,
controle, bombas, baterias, sensores, arquivo e contenção), cada um com
entradas, saídas, pontos de interesse e iluminação. A geração é determinística
por semente --- a mesma semente produz sempre o mesmo laboratório ---, e a
conectividade é garantida por construção e verificada por busca em largura:
nenhuma sala fica inalcançável (Figura~\ref{fig:gerador}). Essa combinação de
determinismo e garantia de conectividade é o que torna a geração testável
automaticamente, e não apenas plausível.

\begin{figure}[htbp]
\centering
\includegraphics[width=0.72\textwidth]{06_gerador.png}
\caption{Geração procedural: semente, módulos, verificação de conectividade por busca em largura e layout final.}
\label{fig:gerador}
\end{figure}

\section{\poc{6} --- Recorte vertical}
\label{sec:poc6}

O sexto incremento fechou o arco jogável: tutorial, persistência, áudio e
epílogo.

\subsection{Tutorial e progressão}

O jogo começa com um tutorial de controles apresentado em etapas curtas, que
apresenta ao jogador a movimentação, a câmera e a interação antes que a
missão exija qualquer procedimento. A estrutura dramática em quatro atos
(Figura~\ref{fig:narrativa}) conduz a ficção: a falha inicial, a percepção de
que há algo inconsistente, a revelação sobre a IA e a escolha final.

\begin{figure}[htbp]
\centering
\includegraphics[width=0.95\textwidth]{09_narrativa.png}
\caption{Estrutura dramática em quatro atos, do restabelecimento da energia à escolha final.}
\label{fig:narrativa}
\end{figure}

\subsection{Persistência}

O progresso é salvo em arquivo local e restaurado com fidelidade: o
salvamento inclui o ponto da missão, as \emph{flags} de estado, os estados
dos equipamentos e o conhecimento narrativo acumulado. Após o carregamento, o
mundo é reconciliado --- luzes, portas e equipamentos refletem o estado
restaurado ---, o que é verificado automaticamente.

\subsection{Áudio}

Os efeitos sonoros são originais, gerados por um script do próprio projeto:
coleta, instalação, porta, erro, alarme, interface e ambiência. A ambiência é
um laço contínuo de baixa intensidade, e os efeitos pontuais usam um conjunto
de canais reutilizáveis. Um defeito encontrado e corrigido nesse subsistema
merece registro: como os arquivos de áudio são importados pelo motor em
formato comprimido, o cálculo do ponto de repetição do laço precisava usar a
duração real do arquivo, e não o tamanho do bloco de dados comprimido --- erro
que causava um laço de menos de um segundo em um trecho de quatro segundos. A
correção foi acompanhada de uma verificação de regressão na suíte.

\subsection{Epílogo}

Ao final, o jogador escolhe entre seguir as instruções da IA, questioná-la ou
reiniciar seu núcleo. A escolha é registrada no estado persistente e altera o
texto final exibido. A decisão de projeto --- declarada no escopo do trabalho
--- foi manter a ramificação restrita ao epílogo, em vez de multiplicar
caminhos.

\section{Engenharia de desempenho}
\label{sec:desempenho_dev}

O alvo de execução em computador sem GPU dedicada guiou decisões em todas as
etapas. No motor, adotou-se o renderizador compatível
(\texttt{gl\_compatibility}), geometria simples, número controlado de luzes e
ausência de efeitos de pós-processamento custosos. No asset, o emissivo foi
dosado (luzes de ambiente em intensidade 1,5--1,6; telas em 1,1--1,3) para
não estourar em branco na ausência de \emph{bloom}. No detalhamento, os
efeitos de animação que o formato glTF não suporta (emissão intermitente)
foram convertidos em lógica de script de baixo custo. E a iluminação do nível
segue as posições das luminárias efetivamente modeladas --- a luz nasce onde
o jogador vê a fonte ---, o que evita tanto o gasto desnecessário quanto o
resultado visual incoerente.

O orçamento de tempo por quadro a 60 Hz é de
\begin{equation}
\label{eq:orcamento}
t_{\text{quadro}} = \frac{1}{f} = \frac{1}{60} \approx 16{,}67\ \text{ms},
\end{equation}
\noindent e foi esse orçamento --- considerando lógica de \emph{gameplay},
física simplificada, iluminação em camadas e interface --- que orientou duas
decisões estruturais: manter efeitos de pós-processamento fora do projeto e
converter animações de material em lógica de script barata. A iluminação em
camadas é um compromisso explícito: em vez de muitas luzes dinâmicas com
sombras (custo alto em GPU integrada), usa-se um conjunto reduzido de luzes
sem sombras, compensado por emissivos dosados no asset e por cores de acento
por zona.

\section{Interface e direção visual}
\label{sec:interface}

A interface do jogo foi tratada como instrumento, não como menu genérico de
videogame. O HUD mantém, no canto superior, o relógio da instalação, a
identificação do setor e o estado de energia; abaixo, à esquerda, a lista de
objetivos com marcação de passo concluído --- o mesmo encadeamento da missão
descrita na \autoref{sec:interacao}; e, na base, os convites de interação e de
status. O resultado é um painel que informa sem competir com o cenário.

O terminal da IA recebeu tratamento visual próprio, com cabeçalho de
identificação do núcleo, histórico de mensagens e linha de entrada. A escolha
de exibir a conversa como terminal --- em vez de balões de diálogo --- reforça
a ficção de sistema operacional e acomoda naturalmente a latência da IA:
enquanto a resposta não chega, o cursor permanece ativo, sem quebrar a
imersão.

O inventário também é diegético: em vez de uma grade abstrata, o jogador vê
o kit técnico com seus itens (fusível, módulo e chave), em uma lista curta que
reforça a correspondência entre o que se carrega e o que se instala nos
painéis. A paleta geral é escura, com cores de função --- branco para
iluminação geral, ciano para interfaces da IA, âmbar para máquinas e vermelho
para emergência e travamentos ---, mantendo coerência com a iluminação em
camadas descrita na \autoref{sec:poc4}.

\section{Level design do recorte vertical}
\label{sec:level_design}

O recorte vertical organiza o espaço em torno de um corredor central
(Figura~\ref{fig:mapa}). A sala principal --- onde começa o jogo --- concentra
a bancada, o painel, o CLP, a bomba e a estação da IA; as demais salas são
alcançadas pelo corredor, que funciona como rota principal e como ponto de
orientação. Os pontos de interesse (arquivo, energia, bombas e controle)
ficam nas extremidades, e o percurso final conduz ao núcleo da IA.

\begin{figure}[htbp]
\centering
\begin{tikzpicture}[font=\scriptsize,
  sala/.style={draw, rounded corners=2pt, align=center, inner sep=3pt,
               minimum width=2.5cm, minimum height=0.95cm, fill=black!4},
  seta/.style={-{Stealth[length=1.8mm]}, semithick}]
\node[sala] (cor) {Corredor\\(rota principal)};
\node[sala, left=14mm of cor] (bom) {Bombas (PI)};
\node[sala, right=14mm of cor] (ctl) {Controle --- IA (PI)};
\node[sala, above=10mm of cor] (arq) {Arquivo (PI)};
\node[sala, below=10mm of cor] (ene) {Energia (PI)};
\draw[seta] (bom) -- (cor);
\draw[seta] (cor) -- (ctl);
\draw[seta] (cor) -- (arq);
\draw[seta] (cor) -- (ene);
\draw[seta] (ene) -- (ctl);
\end{tikzpicture}
\caption{Organização espacial do recorte vertical: corredor central, pontos de interesse (PI) nas extremidades e percurso final até o núcleo da IA.}
\label{fig:mapa}
\par\vspace{2pt}{\footnotesize Fonte: elaborado pelo autor (2026), a partir do level design do projeto.}
\end{figure}

Três princípios guiaram o traçado: o jogador deve sempre saber onde está
(landmarks visuais, letreiros e iluminação por zona); deve saber qual é o
objetivo corrente (HUD persistente); e deve reconhecer a interação possível
(convite contextual). O que o jogo deliberadamente \emph{não} entrega de
imediato é a motivação da IA --- a ambiguidade é conteúdo narrativo, não
falha de comunicação.

\section{Persistência: estrutura do salvamento}
\label{sec:persistencia}

O salvamento é um arquivo de texto no diretório do usuário, com estrutura
versionada em campos independentes (Quadro~\ref{quad:save}). A separação
segue as camadas de memória descritas na \autoref{sec:memoria}: o que é
estado persistente entra no arquivo; o que é sessão não entra. O carregamento
reconcilia o mundo --- luzes, portas e equipamentos assumem o estado salvo
---, e essa reconciliação é verificada automaticamente pelo teste da etapa
correspondente.

\begin{quadro}[htbp]
\centering
\caption{Estrutura do arquivo de salvamento (campos principais).}
\label{quad:save}
\begin{verbatim}
{
  "versao": 1,
  "missao": {"passo_atual": 5, "concluida": true},
  "flags": {"energia_restaurada": true, ...},
  "equipamentos": {"bomba": "RUNNING", "clp": "RUNNING", ...},
  "memoria_narrativa": ["bomba", "revelacao", ...],
  "epilogo": "questionar"
}
\end{verbatim}
\par\vspace{2pt}{\footnotesize Fonte: elaborado pelo autor (2026), a partir do sistema de salvamento do projeto.}
\end{quadro}

\section{Generação procedural: algoritmo}
\label{sec:algoritmo_gerador}

O gerador da ala procedural trabalha em duas fases (Quadro~\ref{quad:gerador}):
primeiro constrói uma árvore de expansão a partir da semente, ligando módulos
por portas compatíveis; em seguida, verifica por busca em largura que todos os
módulos são alcançáveis a partir da entrada e que os pontos de interesse
estão em regiões acessíveis. Se a verificação falha, o gerador descarta o
layout e tenta de novo --- com a mesma semente, o resultado é sempre o mesmo,
o que torna a geração testável e os defeitos reproduzíveis.

\begin{quadro}[htbp]
\centering
\caption{Pseudocódigo do gerador procedural (duas fases).}
\label{quad:gerador}
\begin{verbatim}
gerar(seed):
    rng = aleatorio(seed)
    arvore = expandir_modulos(rng)      # portas compativeis
    layout = posicionar(arvore)
    enquanto nao conectado(layout):
        layout = tentar_novamente(rng)  # rng preserva o determinismo
    retornar layout
\end{verbatim}
\par\vspace{2pt}{\footnotesize Fonte: elaborado pelo autor (2026), a partir do gerador do projeto.}
\end{quadro}

\section{Operação: execução, registros e encerramento}
\label{sec:operacao}

O projeto oferece dois scripts de operação: um para iniciar o jogo (e,
quando disponível, o serviço de IA) e outro para encerrar ambos com segurança.
Os registros de execução e os identificadores de processo ficam em um
diretório de trabalho; ao encerrar, o script pede a saída por um arquivo-sinal
que o próprio jogo observa, em vez de enviar sinais do sistema operacional.
Essa decisão --- incomum, mas necessária --- responde a uma restrição do
ambiente: o motor instalado por pacote confinado recusa sinais externos, de
modo que tanto o encerramento assistido quanto a própria suíte de testes
dependem do canal cooperativo. O detalhe é registrado aqui porque afeta
diretamente a reprodutibilidade das execuções.

\section{Subsistema de áudio}
\label{sec:audio_dev}

O áudio é inteiramente original e gerado por um script do projeto, sem
conteúdo de terceiros: sete efeitos que cobrem coleta, instalação, porta,
erro, alarme, interface e ambiência (Quadro~\ref{quad:audio}). A ambiência é
um laço contínuo de baixa intensidade, e os efeitos pontuais compartilham um
conjunto de canais reutilizáveis, evitando sobreposições quando o jogador
interage rapidamente.

\begin{quadro}[htbp]
\centering
\caption{Catálogo de sons do jogo (gerados por script próprio).}
\label{quad:audio}
\footnotesize
\begin{tabular}{@{}>{\raggedright\arraybackslash}p{2.2cm} >{\raggedright\arraybackslash}p{5.4cm} >{\raggedright\arraybackslash}p{5.2cm}@{}}
\toprule
Som & Uso & Característica \\
\midrule
\texttt{pickup} & coleta de item & curto, ascendente \\
\texttt{install} & instalação no painel & curto, mecânico \\
\texttt{door} & abertura da porta & ruído descendente \\
\texttt{error} & ação inválida & tom descendente seco \\
\texttt{alarm} & alerta do sistema & pulsado, mais longo \\
\texttt{ui} & interface e terminal & clique discreto \\
\texttt{hum} & ambiência contínua & laço de baixa frequência \\
\bottomrule
\end{tabular}
\par\vspace{2pt}{\footnotesize Fonte: elaborado pelo autor (2026), a partir do gerador de áudio do projeto.}
\end{quadro}

\section{Considerações do capítulo}
\label{sec:consideracoes_desenvolvimento}

O desenvolvimento seguiu uma trajetória de complexidade crescente: entrada e
movimento; asset e interação; IA com fronteira de segurança; mundo reativo;
geração procedural; e fechamento narrativo. Cada etapa preservou a
possibilidade de executar e demonstrar o jogo completo, o que permitiu que as
evidências fossem coletadas continuamente --- e não apenas ao final. Os
resultados quantitativos dessa construção são apresentados no
\autoref{cap:resultados}.
